% ===========================================================================
\section{Introduction}\label{sec:intro}
% ===========================================================================

\subsection{The sources}
Seven documents are covered.
\begin{itemize}[leftmargin=2em]
\item \cite{MathMag} (1968): an elementary number-theoretic note settling a conjecture of Sutcliffe
  on integers that are multiplied by an integer when their digits are reversed. Section~\ref{sec:arith}.
\item \cite{JMM} (1965), known to us only through its review. Its theorems concern boundary functions of
  homeomorphisms, continuous, Baire and measurable functions in a disk. Theorems~2--6 of \cite{JMM}
  reappear, with proofs, in the dissertation \cite{Diss} (which says so explicitly); Theorem~1 (on
  homeomorphisms) does not, and is proved here in Section~\ref{sec:bf} by the crosscut methods of
  \cite{MMJ,Diss}.
\item \cite{MMJ} (1966): a short proof of McMillan's theorem that the set of curvilinear convergence of
  a continuous function is of type $\Fsd$, and a generalization to cluster sets. Section~\ref{sec:cc}.
\item \cite{Diss} (1967): the dissertation. Its Introduction (Baire functions) is in
  Section~\ref{sec:baire}; Chapter~I in Sections~\ref{sec:cc}, \ref{sec:bf}, \ref{sec:construct};
  Chapter~II in Section~\ref{sec:discont}.
\item \cite{TAMSb} (1969): every $\Fsd$ set is the set of curvilinear convergence of a continuous
  function with a prescribed boundary function (McMillan's problem); appendix on Baire functions into
  the sphere. Sections~\ref{sec:construct} and~\ref{sec:baire}.
\item \cite{TAMSa} (1969): boundary functions of bounded harmonic functions. Section~\ref{sec:harmonic}.
\item \cite{PAMS} (1969): a bounded continuous function in a cube whose set of curvilinear
  convergence is not Borel. Section~\ref{sec:cube}.
\end{itemize}
A concordance between the numbered results of the sources and the items of this document is given in
Appendix~\ref{app:concordance}.

\subsection{How the document is organised}
Several results occur in more than one source in slightly different settings (the disk $\D$, the upper
half-plane $\Hp$, the strip-like region $Q$). We prove each result once, in the most convenient model,
and transfer it to the others by an explicit homeomorphism (Lemma~\ref{lem:transfer}). When two
sources prove the same statement by different arguments, we give the argument that is shorter or more
general, and say so in the concordance.

\emph{Facts.} Results from the literature that the sources use without proof are stated with their
provenance and proved here (Tietze's arc theorem, the crosscut theorem, the Poisson integral, Banach's
theorem, an analytic non-Borel set, Bagemihl's theorem for continuous functions). Two deep theorems that
the sources rely on, the Jordan curve theorem with accessibility and Kierst's theorem on asymptotic
values, are avoided: we replace the arguments that use them by direct ones (Lemmas~\ref{lem:ha-access}
and~\ref{lem:cube-lemma-elementary}), and record the original arguments in remarks.

\emph{Departures from the sources.} We record here the places where our proof differs in substance from
the source.
\begin{enumerate}[label=(\roman*)]
\item The step (B)(III) of \cite{TAMSb} cites a theorem of Bagemihl and McMillan; we use instead the
  approximation theorem of the dissertation (Theorem~\ref{thm:baire1-char}), as \cite{TAMSb} itself
  suggests.
\item The proof of (F)(I) in \cite{TAMSb} cites an invariance-of-domain corollary from \cite{Newman}; we
  write down the inverse maps explicitly (Lemma~\ref{lem:fold}).
\item The extension theorems for Baire-one functions (cited from \cite{Hausdorff} in \cite{Diss} and
  used in Appendix~(b),(c) of \cite{TAMSb}) are proved (Theorem~\ref{thm:extension}).
\item The ``Figure~5'' argument of \cite{Diss} and the monotonicity of first hitting points are derived
  from a single imported crosscut theorem (Theorem~\ref{fact:crosscut}).
\item Sutcliffe's Theorems~2 and~3, used in \cite{MathMag}, are proved with explicit solutions
  (Propositions~\ref{prop:sut-two-composite}, \ref{prop:sut-explicit}).
\item In the existence of special functions (statement (XX) of \cite{TAMSb}) we give a non-recursive
  construction: all disjointness constraints turn out to be independent of the other member's size.
\end{enumerate}

\subsection{Conventions}\label{sec:conventions}
$\N=\{1,2,3,\dots\}$. ``Countable'' includes ``finite''. We identify $\R^2$ with $\C$ and write points as
$z=x+iy=\langle x,y\rangle$. We use
\begin{align*}
\D&=\{z:\abs z<1\}, & C&=\{z:\abs z=1\}, & \Hp&=\{z:\operatorname{Im}z>0\},\\
Q&=\{\langle x,y\rangle:-1\le x\le1,\ 0<y\le1\}, & Q^\circ&=(-1,1)\times(0,1), &
X&=(-1,1)\times\{0\}.
\end{align*}
The real axis is identified with $\R$; in Sections~\ref{sec:construct} we identify $X$ with $(-1,1)$.
For $n\in\N$, $\Hp_n=\{0<y<1/n\}$, $\overline{\Hp}_n=\{0\le y\le 1/n\}$ and $X_n=\{y=1/n\}$.
$\Sph=\{v\in\R^3:\abs v=1\}$ is the Riemann sphere with the Euclidean metric.
$S(r,p)$ is the open ball of radius $r$ about $p$; $\cl E$ or $\overline E$ is closure, $E^*$ or
$\inter E$ interior (for subsets of $C$ or of $\R$, relative to $C$ or $\R$).
For an interval or arc $I$, $I^*$ is its relative interior; an interval or arc is
\emph{nondegenerate} if it has more than one point.
